\section{Machine-Learning Models and Validation}
\label{sec:models}

\subsection{Thesis protocol: random split}
In the original thesis work, four regressors were trained on the thesis set with the thesis descriptors (248 feature columns: temperature, pressure, 10 DFT descriptors and 236 atom-count columns):
\begin{itemize}
    \item \textbf{Ridge regression}\cite{hoerl1970ridge} with $\alpha = 1$, as a linear baseline;
    \item \textbf{Random Forest}\cite{breiman2001random} with 100 trees;
    \item \textbf{Gradient Boosting}\cite{friedman2001greedy} with 100 trees;
    \item \textbf{a neural network}: a fully connected network with hidden layers of 128, 64 and 32 ReLU units and dropout of 0.2 after the second layer, trained with Adam\cite{kingma2015adam} on mean squared error for 100 epochs with a batch size of 32, in Keras\cite{chollet2015keras}.
\end{itemize}
The first three use scikit-learn\cite{pedregosa2011scikit} with default settings apart from those listed. The dataset was split at random into 80\% training (9,004 points) and 20\% test (2,252 points) with a fixed seed, and features were standardised with a scaler fitted on the training set (\path{notebooks/thesis/17-final-model-comparison.ipynb}). These results are reported for comparison with the thesis only, because a random split of points does not measure prediction for new ILs (Section~\ref{sec:grouped}).

\subsection{Main protocol: grouped cross-validation}
\label{sec:grouped-method}
A typical IL contributes tens to hundreds of points along a few isotherms. With a random split, almost every test point has neighbours from the same IL, and often the same isotherm, in the training set, so the test score measures interpolation along a known isotherm. The main results therefore use the curated dataset (Section~\ref{sec:curated}) and four 5-fold cross-validation schemes that differ in what is held out together\cite{meredig2018can, cawley2010overfitting}:
\begin{itemize}
    \item \textbf{point-wise}: random points, as in the thesis protocol;
    \item \textbf{leave-IL-out}: all points of an IL;
    \item \textbf{leave-cation-out}: all ILs that share a cation;
    \item \textbf{leave-anion-out}: all ILs that share an anion.
\end{itemize}
Leave-anion-out is the most demanding test and the most relevant one, because the anion controls physical \co{} solubility most strongly\cite{cadena2004why, anthony2005anion} and a screening study would vary the anion. The folds are identical for every feature set and model, so differences between feature sets are not due to different splits. $R^2$ and the mean absolute error (MAE, in mole fraction) are reported as the mean and standard deviation over the five folds.

Two tree ensembles were used: a Random Forest with 300 trees and a histogram-based gradient-boosting model\cite{ke2017lightgbm} (scikit-learn \texttt{HistGradientBoostingRegressor}, 500 iterations, learning rate 0.05). Six feature sets were compared, all with temperature and pressure ($T$, $P$):
\begin{itemize}
    \item $T$, $P$ only;
    \item $T$, $P$ and the 15 varying atom counts, counted from the optimised ion structures (ion identity and size, no quantum chemistry);
    \item $T$, $P$ and the 10 thesis DFT descriptors;
    \item $T$, $P$ and the 10 corrected DFT descriptors;
    \item each of the two DFT sets together with the atom counts.
\end{itemize}
The analysis is in \path{scripts/compare_descriptor_sets.py}; per-anion errors of the leave-anion-out predictions are stored with the results in \path{paper/data/}.
